\documentclass[12pt]{article}
\usepackage[utf8]{inputenc}
\usepackage[T1]{fontenc}
\usepackage[margin=1in]{geometry}
\usepackage{setspace}
\usepackage{xcolor}
\usepackage[colorlinks=true, urlcolor=blue!60!black, linkcolor=blue!60!black]{hyperref}
\setlength{\parindent}{0em}
\setlength{\parskip}{0.45em}
\renewcommand{\baselinestretch}{1.0}
\begin{document}
\begin{center}
{\Large \textit{Miller v. California}, 413 U.S. 15 (1973)\par}
\vspace{0.3em}
{\normalsize Case Brief --- One-Page Edition\par}
\end{center}

\textbf{Memory Jogger.} The modern obscenity test: prurient interest, patent offensiveness, and no serious value, under local community standards (\href{https://supreme.justia.com/cases/federal/us/413/15/}{413 U.S. 15}).

\textbf{Facts.} Marvin Miller mailed unsolicited brochures of explicit sexual pictures to a Newport Beach restaurant and was convicted of distributing obscene matter under California law.

\textbf{Procedural History.} The jury applied statewide standards; the Appellate Department summarily affirmed.

\textbf{Votes.} Vacated and remanded, 5--4: Burger for the Court; Douglas and Brennan (with Stewart and Marshall) dissented.

\textbf{Holding.} Obscenity is unprotected, judged by the three-part test below for sexual conduct specifically defined by state law; the \textit{Memoirs} ``utterly without redeeming social value'' test is rejected, and no national standard is required.

\textbf{Reasoning.} Burger sought to end the Court's role as ``an unreviewable board of censorship for the 50 States''; \textit{Memoirs} forced prosecutors to prove a negative, and a national standard would be ``an exercise in futility.'' Douglas found no obscenity exception in the text; Brennan concluded that no formulation could be precise enough.

\textbf{Examples.} \emph{Same-Side:} unsolicited hard-core pamphlets. \emph{Opposite-Side:} a graphic medical textbook. \emph{Fence-Sitter:} an explicit online film with a political monologue.

\textbf{Critique.} Libertarians fault vagueness and chilling, worse online; MacKinnon says it targets offense, not harm to women.

\textbf{Key Quotation.} ``The basic guidelines for the trier of fact must be: (a) whether `the average person, applying contemporary community standards' would find that the work, taken as a whole, appeals to the prurient interest, \textit{Kois v. Wisconsin}, supra, at 230, quoting \textit{Roth v. United States}, supra, at 489; (b) whether the work depicts or describes, in a patently offensive way, sexual conduct specifically defined by the applicable state law; and (c) whether the work, taken as a whole, lacks serious literary, artistic, political, or scientific value.'' --- Chief Justice Burger, Opinion of the Court, 413 U.S. at 24.

\end{document}
